\uuid{Jw2M}
\titre{Points critiques et extremums locaux de $f(x,y)=xy(x+y-1)$}
\niveau{L2}
\module{Analyse}
\chapitre{Fonctions de plusieurs variables}
\sousChapitre{Optimisation sans contrainte}
\theme{Points critiques, matrice hessienne, points selles, minimum local}
\auteur{Maxime Nguyen}
\datecreate{2026-03-24}
\organisation{amscc}
\difficulte{3}

\contenu{
\texte{
  On considère la fonction $f : \mathbb{R}^2 \to \mathbb{R}$ définie par $f(x,y) = xy(x+y-1)$.
}
\begin{enumerate}
  \item \question{Déterminer les points critiques de $f$. Admet-elle des extremums locaux ?}
    \reponse{
      On calcule les dérivées partielles :
      \[
        \frac{\partial f}{\partial x}(x,y) = y(2x+y-1), \qquad \frac{\partial f}{\partial y}(x,y) = x(2y+x-1).
      \]
      On résout $\nabla f = 0$ :
      \begin{itemize}
        \item $x = 0 \Rightarrow y = 0$ ou $y = 1$ ;
        \item $y = 0 \Rightarrow x = 0$ ou $x = 1$ ;
        \item $x \neq 0$ et $y \neq 0$ : le système $2x+y-1=0$, $x+2y-1=0$ donne $x = y = \tfrac{1}{3}$.
      \end{itemize}
      Les points critiques sont $(0,0)$, $(0,1)$, $(1,0)$ et $\bigl(\tfrac{1}{3}, \tfrac{1}{3}\bigr)$.

      Les dérivées secondes sont :
      \[
        \frac{\partial^2 f}{\partial x^2} = 2y, \quad \frac{\partial^2 f}{\partial y^2} = 2x, \quad \frac{\partial^2 f}{\partial x \partial y} = 2x+2y-1.
      \]
      \begin{itemize}
        \item En $(0,0)$ : $\Delta_f(0,0) = -1 < 0$ $\Rightarrow$ point selle.
        \item En $(0,1)$ : $\Delta_f(0,1) = -1 < 0$ $\Rightarrow$ point selle.
        \item En $(1,0)$ : $\Delta_f(1,0) = -1 < 0$ $\Rightarrow$ point selle.
        \item En $\bigl(\tfrac{1}{3},\tfrac{1}{3}\bigr)$ : $\Delta_f = \tfrac{1}{3} > 0$ et $\dfrac{\partial^2 f}{\partial x^2} = \tfrac{2}{3} > 0$ $\Rightarrow$ minimum local.
      \end{itemize}
      Conclusion : $(0,0)$, $(0,1)$, $(1,0)$ sont des points selles et $\bigl(\tfrac{1}{3},\tfrac{1}{3}\bigr)$ est un minimum local.
    }
\end{enumerate}
}
